\section{Installation}\label{sec:installation}

PyRATE-TA requires \textbf{Python~3.12 or later}. It can be installed directly from \href{https://pypi.org/project/pyrate-ta/}{PyPI} or installed from source using \href{https://docs.astral.sh/uv/}{\textbf{uv}}.

\subsection{Quick install via PyPI (Recommended)}\label{sec:install-pypi}

To install the latest official release into your active Python environment:

\begin{lstlisting}[language=bash, style=py]
pip install pyrate-ta
# or with uv:
uv pip install pyrate-ta

# Register bundled fonts into matplotlib:
pymorgan-install-fonts
\end{lstlisting}

\noindent You can also launch the graphical interface directly without prior installation using \code{uvx}:

\begin{lstlisting}[language=bash, style=py]
uvx --from pyrate-ta pyrate-ta-gui
\end{lstlisting}

\subsection{Installation from Source (Development)}

From the repository root:

\begin{lstlisting}[language=bash, style=py]
uv venv                          # create .venv (Python >= 3.11)
uv pip install -e ".[gui,dev]"   # editable install + dependencies
uv run pymorgan-install-fonts    # PyMORGAN's fonts, in *this* environment
\end{lstlisting}

The font step is not optional in practice. PyRATE-TA uses PyMORGAN's font stack and
ships no second installer, but Matplotlib registers fonts per environment: skip
it and every figure falls back to DejaVu Sans, which is why a PyRATE-TA plot can
look unlike a PyMORGAN one while both claim the same style. The GUI checks at
start-up and logs the command if the preferred fonts are missing.

uv automatically uses the project's \code{.venv}: prefix one-off commands with
\code{uv run} to avoid activating it manually.

\subsection{Co-development with PyMORGAN}

PyRATE-TA depends on PyMORGAN. While the two are developed side by side under
\code{Scripts/}, PyMORGAN is picked up from the sibling checkout as an editable
path dependency, so an upstream edit is visible immediately without a
reinstall:

\begin{lstlisting}[style=py]
# pyproject.toml
[tool.uv.sources]
pymorgan = { path = "../PyMORGAN", editable = true }
\end{lstlisting}

Replace that block with a plain version pin once PyMORGAN is published to an
index. The same editable dependency puts \code{pymorgan-install-fonts} on
PyRATE-TA's path, which is how the fonts are registered for this environment.

\begin{table}[h]
  \centering
  \begin{tabular}{lll}
    \toprule
    Package & Minimum version & Purpose \\
    \midrule
    \code{pymorgan}   & 0.6    & Loading, datasets, plotting, aesthetics \\
    \code{numpy}      & 2.1.0  & Array operations \\
    \code{scipy}      & 1.14.1 & Optimisation and matrix exponentials \\
    \code{tomlkit}    & 0.13   & Settings file read/write \\
    \code{PyQt6}      & 6.7.1  & Graphical interface framework \\
    \code{PyQt6\_sip} & 13.8.0 & Qt bindings interface \\
    \bottomrule
  \end{tabular}
  \caption{Core dependencies installed by \code{uv pip install -e .}}
  \label{tab:dependencies}
\end{table}

\subsection{Console scripts}

The editable install puts the following entry points on the path:

\begin{table}[h]
  \centering
  \begin{tabular}{ll}
    \toprule
    Command & Purpose \\
    \midrule
    \code{pyrate-ta-gui}      & Launch the fitting interface \\
    \code{pyrate-ta-edit-gui} & Open \code{main\_window.ui} in Qt Designer \\
    \code{pyrate-ta-settings} & Standalone editor for the analysis settings \\
    \bottomrule
  \end{tabular}
  \caption{Console scripts, mirroring PyMORGAN's. Presentation settings are
    edited with \code{pymorgan-settings}, not here.}
  \label{tab:scripts}
\end{table}

\subsection{Verifying the installation}

\begin{lstlisting}[language=bash, style=py]
uv run python -c "import pyrate_ta; print(pyrate_ta.__version__)"
uv run pytest
uv run python scripts/run_checks.py
\end{lstlisting}

\code{run\_checks.py} exercises the settings round-trip, the lazy public API,
the PyMORGAN boundary and the integrity of the Qt Designer layout; the Qt
checks are skipped, not failed, on a headless machine.
\code{python scripts/run\_changed\_tests.py} runs only the suites affected by
the current diff, and falls back to the full suite whenever a change cannot be
mapped --- an unmapped change is never reported as tested.

\subsection{Versioning}

Every edit to the codebase must bump the version in
\code{src/pyrate_ta/\_\_about\_\_.py}, using \code{python scripts/bump\_version.py}
rather than editing it by hand. The scheme is \code{0.x.yymmdd.devN} (PEP~440):
\code{yymmdd} is the build date and \code{N} the iteration within that day,
restarting at 1 on a new date. A bare letter suffix (\code{0.1.260728d}) is not
valid PEP~440 and breaks the build. \code{--check} exits non-zero when the
recorded version was not produced today, which suits a pre-commit hook.
